\documentclass[11pt,a4paper]{article}

% ------------------------------------------------
% Encoding and Fonts
% ------------------------------------------------

\usepackage[utf8]{inputenc}
\usepackage[T1]{fontenc}
\usepackage{lmodern}

% ------------------------------------------------
% Mathematics
% ------------------------------------------------

\usepackage{amsmath}
\usepackage{amssymb}
\usepackage{amsthm}
\usepackage{mathtools}

% ------------------------------------------------
% Page Layout
% ------------------------------------------------

\usepackage[margin=1in]{geometry}
\usepackage{microtype}

% ------------------------------------------------
% Additional Packages
% ------------------------------------------------

\usepackage{hyperref}
\usepackage{enumitem}
\usepackage{graphicx}
\usepackage{booktabs}

% ------------------------------------------------
% Paragraph Formatting
% ------------------------------------------------

\setlength{\parindent}{0pt}
\setlength{\parskip}{6pt}

% ------------------------------------------------
% Theorem Environments
% ------------------------------------------------

\newtheorem{theorem}{Theorem}[section]
\newtheorem{lemma}[theorem]{Lemma}
\newtheorem{proposition}[theorem]{Proposition}
\newtheorem{corollary}[theorem]{Corollary}

\theoremstyle{definition}
\newtheorem{definition}[theorem]{Definition}
\newtheorem{example}[theorem]{Example}

\theoremstyle{remark}
\newtheorem{remark}[theorem]{Remark}

\setlength{\parindent}{0pt}
\setlength{\parskip}{6pt}
% ------------------------------------------------
% Title
% ------------------------------------------------


% --------------------------------------------------
% Title
% --------------------------------------------------
\title{
    \textbf{Optimal Relaxation Parameters for the
    Successive Over-Relaxation Method}\\[4pt]
    \large Theory, Derivation, and Numerical Experiments
}

\author{Biajid Chowdhury}
\date{\today}

\begin{document}

\maketitle

% --------------------------------------------------
% Abstract
% --------------------------------------------------
\begin{abstract}

One of the most fundamental problems in computational
science and engineering is solving a system of linear
equations of the form
\[
    Ax=b.
\]
For nearly two millennia, engineers, mathematicians and scientists
have explored different approaches to solving such
systems efficiently at an acceptable computational cost. From the elimination procedures
of ancient mathematics to modern direct and iterative
algorithms, this endeavor has produced numerous
important developments and continues to motivate
research in scientific computing.

Among these developments is the introduction of
relaxation parameters into iterative methods and,
more importantly, the problem of determining their
optimal values. The successive over-relaxation
(SOR) method provides a particularly interesting
example, since a suitable choice of the relaxation
parameter can substantially accelerate convergence,
whereas a poor choice may lead to extremely slow
convergence or even divergence.

The mathematical foundations of SOR emerged from
several decades of research on iterative methods
and numerical solutions of differential equations.
An important breakthrough came through the work
of David M.~Young, whose doctoral research in 1950
and subsequent publications established a systematic
theoretical framework for successive over-relaxation.
For certain structured classes of matrices, this
theory provides an explicit expression for the
optimal relaxation parameter in terms of the
spectral radius of the Jacobi iteration matrix.
For general matrices, however, determining the
optimal parameter remains a more difficult problem,
and no comparably simple universal formula is
available.

In this note, we study the mathematical formulation
of the Jacobi, Gauss--Seidel, and successive
over-relaxation methods, with particular attention
to their convergence properties. We derive the
optimal SOR parameter for an important class of
tridiagonal matrices arising from finite difference
discretizations of boundary value problems.
Finally, we implement the methods in MATLAB and
compare their numerical performance, examining
the relationship between relaxation parameters,
spectral radii, convergence rates, and computational
accuracy.

The development of optimal relaxation methods
represents a remarkable chapter in the history
of numerical computation. Its mathematical
elegance, practical significance, and long
historical development deserve closer attention.
It is therefore somewhat surprising that the
derivation of optimal relaxation parameters
receives relatively limited treatment in many
undergraduate textbooks, and instructors often overlook it. We hope
that the present discussion will help illuminate
both the mathematics and the significance of
this important subject.

\end{abstract}

\tableofcontents
\newpage



\section{Introduction}

From elementary school, we begin learning how to solve
equations. In fact, the journey probably starts as early
as the fourth grade, when we encounter a simple problem
such as
\[
    4x=8.
\]
By middle school, we are already accustomed to word
problems involving several unknown quantities. We
translate the information given in words into a system
of equations and then try to determine the unknowns.
Without necessarily knowing its name, we are often
using elementary forms of the elimination procedure
that later becomes Gaussian elimination.

By high school, many students have encountered systems
of linear equations, determinants, and perhaps even
Cramer's rule. Thus, the mathematical problem we are
going to investigate in this note is certainly nothing
new to the reader. What changes is not the problem
itself, but its size and the computational resources
required to solve it.

Consider a system of $n$ linear equations in $n$
unknowns, written in matrix form as
\[
    Ax=b,
\]
where $A\in\mathbb{R}^{n\times n}$,
$x\in\mathbb{R}^n$ is the unknown vector, and
$b\in\mathbb{R}^n$ is given.

For a system involving two or three unknowns, Gaussian
elimination is straightforward. Even for a system of
a few hundred equations, a modern computer can solve
the problem with little difficulty. But what happens
when the number of unknowns reaches several million?
Such systems arise naturally in computational science
and engineering, particularly when differential
equations are approximated by finite difference or
finite element methods.

The elimination procedure developed by Gauss and
subsequently extended into the Gauss--Jordan method
remains one of the fundamental achievements of
numerical linear algebra. However, for a general
dense $n\times n$ matrix, Gaussian elimination
requires on the order of $n^3$ arithmetic operations
and $n^2$ storage locations. As $n$ becomes very
large, these requirements can become prohibitively
expensive. Although modern direct methods exploit
sparsity and other special matrix structures to
reduce these costs, it is natural to ask whether
we can approach the problem differently.

What alternative might we consider?

Perhaps we could return to something familiar from
our earlier mathematical education: factorization.
Instead of solving $Ax=b$ directly, suppose we
could write
\[
    A=LU,
\]
where $L$ is lower triangular and $U$ is upper
triangular. Then the original system becomes
\[
    LUx=b,
\]
which can be solved through two triangular systems:
\[
    Ly=b,\qquad Ux=y.
\]

This certainly simplifies the solution once the
factors are known. But obtaining those factors is
itself a computational problem. In fact, Gaussian
elimination is precisely one of the principal
procedures for constructing an $LU$ factorization.
Thus, factorization does not automatically remove
the computational difficulty; rather, it provides
another way of organizing the same fundamental
calculation.

So we ask a different question:

\emph{Must we determine the solution in a finite
sequence of elimination steps, or could we begin
with an approximate solution and improve it
repeatedly until it becomes sufficiently accurate?}

This question leads us to iterative methods.

Instead of attempting to solve $Ax=b$ directly,
an iterative method begins with an initial guess
$x^{(0)}$ and constructs a sequence
\[
    x^{(0)},\ x^{(1)},\ x^{(2)},\ \ldots
\]
that, under suitable conditions, converges to
the exact solution $x^*$:
\[
    \lim_{k\to\infty}x^{(k)}=x^*,
    \qquad Ax^*=b.
\]

The idea is simple: start somewhere, make a
correction, and repeat. The mathematical difficulty
lies in designing a correction that guarantees
convergence and, equally importantly, makes that
convergence sufficiently fast.

Two classical examples are the Jacobi and
Gauss--Seidel methods. Both construct successive
approximations by rearranging the equations of
the original system. Their simplicity makes them
attractive, but their convergence can be painfully
slow, particularly for large systems arising from
the discretization of differential equations.

Can we accelerate the process?

One possibility is to modify the correction at
each iteration by introducing a relaxation
parameter $\omega$. Rather than accepting the
new approximation exactly as computed, we move
toward it by an amount controlled by $\omega$.
This leads to the successive over-relaxation
(SOR) method, a modification of Gauss--Seidel
that can converge substantially faster when
the parameter is chosen appropriately.

But this improvement introduces a new problem:
\emph{how should we choose $\omega$?}

A parameter that works remarkably well for one
matrix may perform poorly for another. Even for
the same matrix, a small change in $\omega$ may
significantly affect the number of iterations
required for convergence. Therefore, the problem
is not merely to introduce a relaxation parameter,
but to understand whether an optimal value exists
and how it can be determined mathematically.

This is the central question of the present note.

We begin by developing the Jacobi and Gauss--Seidel
methods through the matrix splitting
\[
    A=D+L+U,
\]
where $D$ contains the diagonal entries of $A$,
while $L$ and $U$ denote its strictly lower and
strictly upper triangular parts, respectively.
We then introduce the SOR method and derive its
iteration matrix. Using the spectral radius, we
investigate the convergence of these stationary
iterative methods and examine how the relaxation
parameter influences the convergence rate.

Our principal objective is to derive an explicit
expression for the optimal SOR parameter for a
particular class of structured matrices. We
then apply the theory to a tridiagonal system
obtained from the finite difference discretization
of a second-order boundary value problem. In
this setting, the eigenvalues of the Jacobi
iteration matrix can be determined analytically,
allowing us to obtain the optimal relaxation
parameter in closed form.

Finally, we implement the Jacobi, Gauss--Seidel,
and SOR methods in MATLAB and compare their
numerical performance. Particular attention
is given to iteration counts, convergence
rates, and the distinction between iteration
error and discretization error.

Our aim is not to present SOR merely as another
algorithm for solving linear equations. Rather,
we wish to understand the mathematical reasoning
behind its acceleration, the role of eigenvalues
in predicting convergence, and why the seemingly
simple question of choosing a relaxation parameter
leads to such an interesting problem in numerical
linear algebra.

%``latex
\section{Iterative Procedures for Solving Systems of Linear Equations}

Before introducing iterative methods for solving linear systems,
let us begin with a rather elementary problem from computer science.

Suppose we are given the array
\[
    [1,\ 10,\ 2,\ 5,\ 9,\ -1]
\]
and asked to determine whether the number $9$ appears in it.

How should we proceed?

An enthusiastic reader familiar with algorithms might immediately
suggest sorting the array and then performing a binary search.
That is certainly a reasonable approach. But before choosing an
algorithm, perhaps we should ask another question:

\emph{Are we going to perform this search only once, or will we
need to search the same array repeatedly?}

If the search is performed only once, sorting the entire array
may be unnecessary. We can simply begin with the first element,
compare it with $9$, and continue until either the desired
number is found or the end of the array is reached. This is
the familiar linear search algorithm.

However, if we expect to search the same array repeatedly,
perhaps every minute, the situation changes. In that case,
it may be worthwhile to sort the array once and then use
binary search for subsequent queries. By examining the
middle element and repeatedly eliminating half of the
remaining search interval, binary search requires only
$O(\log n)$ comparisons, compared with $O(n)$ comparisons
in the worst case for linear search.

The important lesson is that the choice of an algorithm
depends on the circumstances in which it will be used.
A method that is attractive for one computation may
not be the most efficient choice for another.

A similar consideration arises when solving systems of
linear equations. If the same coefficient matrix is used
with many different right-hand sides, computing a matrix
factorization once and reusing it may be particularly
advantageous. On the other hand, for a very large sparse
system that needs to be solved only once, constructing
and storing such a factorization may be unnecessarily
expensive.

What alternative do we have?

Instead of attempting to obtain the solution through
a direct procedure, we might begin with an approximate
solution and improve it repeatedly. If the approximations
converge sufficiently close to the desired solution
within our computational budget, we accept the result.
Otherwise, we report that the prescribed convergence
criterion has not been satisfied.

This is the fundamental idea behind an iterative method.

Of course, unlike searching an array, we are not examining
a collection of candidate solutions one by one. Rather,
we construct a sequence of approximations according to
a mathematical rule designed to bring us closer to the
solution.

Let us see how this works through a simple example.


\subsection{A Simple Example}

Consider the system
\begin{equation}
\begin{aligned}
    4x+y &= 9,\\
    x+3y &= 7.
\end{aligned}
\label{eq:simple_system}
\end{equation}

In matrix notation,
\[
    \begin{bmatrix}
        4 & 1\\
        1 & 3
    \end{bmatrix}
    \begin{bmatrix}
        x\\y
    \end{bmatrix}
    =
    \begin{bmatrix}
        9\\7
    \end{bmatrix}.
\]

We could solve this system immediately using Gaussian
elimination. In fact, its exact solution is
\[
    \boxed{
        x=\frac{20}{11},
        \qquad
        y=\frac{19}{11}.
    }
\]

But suppose we did not know the solution and wanted
to construct it through successive approximations.

Rearranging the equations in \eqref{eq:simple_system},
we obtain
\begin{equation}
\begin{aligned}
    x &= \frac{9-y}{4},\\
    y &= \frac{7-x}{3}.
\end{aligned}
\label{eq:simple_rearranged}
\end{equation}

These equations suggest an iterative procedure.

Choose an initial approximation
\[
    x^{(0)}=0,
    \qquad
    y^{(0)}=0.
\]

Now use the right-hand sides of
\eqref{eq:simple_rearranged} to construct new
approximations.

For the first iteration,
\[
\begin{aligned}
    x^{(1)}
    &=\frac{9-y^{(0)}}{4}
      =\frac94=2.25,\\
    y^{(1)}
    &=\frac{7-x^{(0)}}{3}
      =\frac73\approx2.333333.
\end{aligned}
\]

For the second iteration,
\[
\begin{aligned}
    x^{(2)}
    &=\frac{9-y^{(1)}}4
      \approx1.666667,\\
    y^{(2)}
    &=\frac{7-x^{(1)}}3
      \approx1.583333.
\end{aligned}
\]

Continuing in this manner gives the following
sequence of approximations:

\begin{center}
\begin{tabular}{ccc}
\toprule
Iteration $k$ & $x^{(k)}$ & $y^{(k)}$\\
\midrule
0 & 0.000000 & 0.000000\\
1 & 2.250000 & 2.333333\\
2 & 1.666667 & 1.583333\\
3 & 1.854167 & 1.777778\\
4 & 1.805556 & 1.715278\\
5 & 1.821181 & 1.731481\\
\midrule
Exact & 1.818182 & 1.727273\\
\bottomrule
\end{tabular}
\end{center}

Notice that the approximations move toward the
exact solution. Although they do not approach it
monotonically, the errors become progressively smaller.

The procedure we have just constructed is an example
of the \emph{Jacobi iterative method}.

But an important question remains:

\emph{Will this procedure always converge, regardless
of the system of equations?}

Unfortunately, the answer is no.

The convergence of an iterative method depends on
the structure of the coefficient matrix and the
particular iteration rule being used. We shall
investigate these conditions mathematically after
developing the general formulation.


\subsection{The Jacobi Iterative Method}

Consider the linear system
\begin{equation}
    Ax=b,
    \label{eq:general_system}
\end{equation}
where $A=(a_{ij})\in\mathbb{R}^{n\times n}$.

Suppose that every diagonal entry of $A$ is nonzero:
\[
    a_{ii}\neq0,
    \qquad i=1,\ldots,n.
\]

The $i$th equation of the system is
\[
    a_{i1}x_1+\cdots+a_{ii}x_i
    +\cdots+a_{in}x_n=b_i.
\]

Solving for $x_i$ gives
\[
    x_i=
    \frac{1}{a_{ii}}
    \left(
        b_i-\sum_{\substack{j=1\\j\neq i}}^n
        a_{ij}x_j
    \right).
\]

This suggests the iteration
\begin{equation}
    \boxed{
    x_i^{(k+1)}
    =
    \frac{1}{a_{ii}}
    \left(
        b_i-
        \sum_{\substack{j=1\\j\neq i}}^n
        a_{ij}x_j^{(k)}
    \right).
    }
    \label{eq:jacobi_component}
\end{equation}

The essential feature of the Jacobi method is that
every component of the new approximation is computed
using only the components of the previous approximation.

Thus, during iteration $k+1$, the values
\[
    x_1^{(k+1)},\ldots,x_n^{(k+1)}
\]
are calculated independently from the vector $x^{(k)}$.

This makes the method conceptually simple and
particularly suitable for parallel computation.

\subsubsection*{Matrix Formulation}

Write
\[
    A=D+L+U,
\]
where $D$ is the diagonal part of $A$, $L$ is its
strictly lower triangular part, and $U$ is its
strictly upper triangular part.

Then
\[
    (D+L+U)x=b,
\]
or equivalently,
\[
    Dx=b-(L+U)x.
\]

Replacing the unknown vector on the right-hand
side by the current approximation gives
\[
    Dx^{(k+1)}
    =
    b-(L+U)x^{(k)}.
\]

Since $D$ is invertible,
\begin{equation}
    \boxed{
    x^{(k+1)}
    =
    -D^{-1}(L+U)x^{(k)}
    +D^{-1}b.
    }
    \label{eq:jacobi_matrix}
\end{equation}

Define
\[
    T_J=-D^{-1}(L+U),
    \qquad
    c_J=D^{-1}b.
\]

The Jacobi iteration therefore takes the compact form
\begin{equation}
    \boxed{
    x^{(k+1)}=T_Jx^{(k)}+c_J.
    }
    \label{eq:jacobi_compact}
\end{equation}

The matrix $T_J$ is called the
\emph{Jacobi iteration matrix}.

As we shall see later, its eigenvalues determine
whether the iteration converges.


\subsection{The Gauss--Seidel Iterative Method}

Let us return briefly to our elementary example.

Recall that the Jacobi method uses the previous
approximation to calculate every component of
the new vector.

But suppose we have already calculated a new
value of $x$. Why should we continue using the
old value when calculating $y$?

Would it not be more reasonable to use the
most recently available information?

This simple observation leads to the
Gauss--Seidel method.

For the system
\[
\begin{aligned}
    4x+y&=9,\\
    x+3y&=7,
\end{aligned}
\]
we now use
\begin{equation}
\begin{aligned}
    x^{(k+1)}
    &=\frac{9-y^{(k)}}4,\\
    y^{(k+1)}
    &=\frac{7-x^{(k+1)}}3.
\end{aligned}
\label{eq:gs_example}
\end{equation}

Notice the difference: the second equation uses
$x^{(k+1)}$, not $x^{(k)}$.

Starting again from
\[
    x^{(0)}=y^{(0)}=0,
\]
the first iteration gives
\[
\begin{aligned}
    x^{(1)}&=\frac94=2.25,\\
    y^{(1)}&=\frac{7-2.25}{3}
             \approx1.583333.
\end{aligned}
\]

The second iteration gives
\[
\begin{aligned}
    x^{(2)}
    &=\frac{9-1.583333}{4}
      \approx1.854167,\\
    y^{(2)}
    &=\frac{7-1.854167}{3}
      \approx1.715278.
\end{aligned}
\]

Continuing the process, we obtain

\begin{center}
\begin{tabular}{ccc}
\toprule
Iteration $k$ & $x^{(k)}$ & $y^{(k)}$\\
\midrule
0 & 0.000000 & 0.000000\\
1 & 2.250000 & 1.583333\\
2 & 1.854167 & 1.715278\\
3 & 1.821181 & 1.726273\\
4 & 1.818432 & 1.727189\\
5 & 1.818203 & 1.727266\\
\midrule
Exact & 1.818182 & 1.727273\\
\bottomrule
\end{tabular}
\end{center}

Comparing these values with the Jacobi iterations,
we observe that Gauss--Seidel converges considerably
faster for this particular system.

This improvement comes from using the newly
computed components immediately rather than
waiting until the next iteration.

\subsubsection*{General Formulation}

For a general system $Ax=b$, the Gauss--Seidel
iteration is
\begin{equation}
    \boxed{
    x_i^{(k+1)}
    =
    \frac{1}{a_{ii}}
    \left(
        b_i
        -
        \sum_{j=1}^{i-1}
        a_{ij}x_j^{(k+1)}
        -
        \sum_{j=i+1}^{n}
        a_{ij}x_j^{(k)}
    \right).
    }
    \label{eq:gs_component}
\end{equation}

The first summation uses newly computed values,
while the second uses values from the previous
iteration.

Using the decomposition
\[
    A=D+L+U,
\]
we can write
\[
    (D+L)x^{(k+1)}
    =
    b-Ux^{(k)}.
\]

Therefore,
\begin{equation}
    \boxed{
    x^{(k+1)}
    =
    -(D+L)^{-1}Ux^{(k)}
    +(D+L)^{-1}b.
    }
    \label{eq:gs_matrix}
\end{equation}

Define
\[
    T_{GS}=-(D+L)^{-1}U,
    \qquad
    c_{GS}=(D+L)^{-1}b.
\]

Then
\begin{equation}
    \boxed{
    x^{(k+1)}
    =
    T_{GS}x^{(k)}+c_{GS}.
    }
    \label{eq:gs_compact}
\end{equation}

Although the formula contains $(D+L)^{-1}$,
we do not explicitly compute this inverse in
practice. Instead, each iteration requires
solving a lower triangular system, which can
be accomplished efficiently through forward
substitution.


\subsection{A Common Mathematical Structure}

The Jacobi and Gauss--Seidel methods may appear
different in their componentwise formulations,
but they share the same fundamental structure:
\[
    \boxed{
    x^{(k+1)}=Tx^{(k)}+c.
    }
\]

Here, $T$ is the iteration matrix and $c$ is
a fixed vector.

Suppose the iteration converges to a vector
$x^*$. Taking limits gives
\[
    x^*=Tx^*+c.
\]

Therefore,
\[
    (I-T)x^*=c.
\]

For both Jacobi and Gauss--Seidel, this fixed-point
equation is equivalent to the original system
$Ax=b$.

But convergence is not automatic.

Let
\[
    e^{(k)}=x^{(k)}-x^*
\]
denote the error at iteration $k$.

Since
\[
    x^{(k+1)}=Tx^{(k)}+c
\]
and
\[
    x^*=Tx^*+c,
\]
subtracting the two equations gives
\[
    e^{(k+1)}=Te^{(k)}.
\]

Repeated application yields
\begin{equation}
    \boxed{
    e^{(k)}=T^ke^{(0)}.
    }
    \label{eq:error_propagation}
\end{equation}

Thus, the convergence problem reduces to
understanding the behavior of the matrix powers
\[
    T,\ T^2,\ T^3,\ldots.
\]

A fundamental result from numerical linear algebra
states that
\[
    T^k\longrightarrow0
    \quad\text{as }k\to\infty
\]
if and only if
\[
    \boxed{\rho(T)<1,}
\]
where $\rho(T)$ denotes the spectral radius of $T$:
\[
    \rho(T)=\max_i|\lambda_i(T)|.
\]

Consequently, a stationary iteration converges
to the unique solution for every initial guess
if and only if the spectral radius of its
iteration matrix is strictly less than one.

This provides the mathematical foundation for
studying and comparing iterative methods.

However, knowing that an iteration converges
is only part of the problem. In computational
science, we are equally interested in how
quickly it converges.

The Gauss--Seidel method improves upon Jacobi
for many important classes of matrices, but
its convergence can still be extremely slow.

Can we modify Gauss--Seidel to obtain a faster
method without fundamentally changing its
simple computational structure?

The answer leads us to the successive
over-relaxation method, which we investigate
in the next section.


\subsection*{A Connection to Optimization and Physical Equilibrium}

For the more advanced reader who may still be wondering
why such iterative procedures should work in the first
place, there is another interesting way to look at the
problem. Suppose that $A$ is a symmetric positive
definite matrix. Then solving
\[
    Ax=b
\]
is equivalent to minimizing the quadratic functional
\[
    \boxed{
    E(x)=\frac12 x^T A x-x^T b.
    }
\]
Indeed, its gradient is
\[
    \nabla E(x)=Ax-b,
\]
and therefore the stationary point satisfies
\[
    \nabla E(x)=0
    \quad\Longleftrightarrow\quad
    Ax=b.
\]
Moreover, since the Hessian is
\[
    \nabla^2 E(x)=A,
\]
and $A$ is positive definite, this stationary point
is the unique global minimum.

This interpretation is closely related to the
principle of minimum potential energy in mechanics.
Many stable physical systems, when displaced from
equilibrium, tend to evolve toward a state of lower
energy. In a certain sense, nature is already
performing its own optimization! We may think of
an iterative numerical method in a similar spirit:
starting from an initial approximation, we repeatedly
correct the current state in the hope of approaching
the desired equilibrium. Some methods make rapid
progress, others move painfully slowly, and some
may fail to converge altogether. Of course, not
every iterative method follows a strictly decreasing
energy path, and the analogy should not be taken
as a general convergence proof. Nevertheless, it
provides a useful intuition for the problem we
are studying. Our task is not merely to move the
approximation toward the solution, but to design
a procedure that makes this journey as smooth,
reliable, and swift as possible. As we shall see,
the relaxation parameter in the SOR method gives
us precisely one way to influence that journey.


One might object that our energy
interpretation applies only when $A$ is symmetric
positive definite. That is certainly true. However,
there is some good news: many important systems
arising from partial differential equations,
particularly those obtained by discretizing
self-adjoint elliptic operators, naturally possess
this structure. For a general matrix, the situation
is somewhat more complicated, but we are not
entirely without alternatives. For example, if
$A$ is nonsingular, we may transform $Ax=b$ into
the equivalent symmetric positive definite system
\[
    A^T A x=A^T b.
\]
We may also consider suitable regularization
techniques when the problem is ill-conditioned
or requires additional stabilization. Of course,
these transformations are not without computational
consequences: forming $A^T A$ squares the
condition number in the Euclidean norm and may
destroy useful sparsity, while regularization
generally modifies the original problem.
Nevertheless, the essential idea remains
valuable. Even when the original system does
not admit a simple minimum-energy interpretation,
optimization and matrix transformations often
provide alternative ways of studying and solving
it. We shall not pursue these generalizations
here, since our principal interest lies in the
classical SOR method and the structured matrices
for which its optimal relaxation parameter can
be determined analytically.


\section{The Successive Over-Relaxation Method}

Before introducing the successive over-relaxation
method mathematically, let us return to a simple
geometric picture.

Suppose someone is standing at the origin
$(0,0)$ in the Cartesian plane and wishes to
reach the point $(1,1)$. They may begin by
choosing a direction, taking a step, examining
their new position, and then calculating the
next step. By repeating this procedure
systematically, they hope eventually to reach
their destination.

This is roughly the geometric intuition behind
the Jacobi and Gauss--Seidel methods. We begin
with an initial approximation, calculate a
new approximation according to a prescribed
rule, and repeat the process in the hope of
approaching the exact solution.

But must we always accept the newly calculated
position exactly as it is?

Suppose our current position is represented
by the vector $x^{(k)}$, and an iterative
procedure suggests a new position
$\widetilde{x}^{(k+1)}$. Instead of moving
directly to that position, we might decide
to travel only $80\%$ of the suggested
displacement. In that case, our next position
would be
\[
    x^{(k+1)}
    =
    0.2x^{(k)}
    +
    0.8\widetilde{x}^{(k+1)}.
\]

In other words, we retain $20\%$ of the
previous approximation and accept $80\%$
of the newly calculated one.

There is a familiar intuition behind this
idea. The reader might remember going on
a trek during their middle school years. When
following an unfamiliar trail, one does
not always rush toward the next visible
point. It is often useful to look ahead
while also keeping track of the path
already traveled. Moving too cautiously
may slow the journey, but moving too
aggressively may lead us away from the
intended route.

A similar consideration arises in
iterative numerical methods. We have
a current approximation and a newly
calculated correction. How much of that
correction should we accept?

Sometimes taking a small corrective step improves
stability. Sometimes it merely slows
convergence.
And, perhaps surprisingly, sometimes moving
\emph{beyond} the newly calculated approximation,
while still retaining a contribution from the
previous approximation, can bring us to the
solution considerably faster. This is precisely
the intuition behind the successive
over-relaxation (SOR) method.

Of course, the trekking analogy is only
an intuition. An iterative method does
not possess a physical sense of direction,
and moving farther does not necessarily
mean moving closer to the solution.
Whether a particular choice accelerates
convergence must ultimately be determined
through mathematical analysis.


\subsection{Introducing the Relaxation Parameter}

Suppose that at iteration $k$ we have
the current approximation $x^{(k)}$,
and a given iterative procedure produces
a candidate approximation
$\widetilde{x}^{(k+1)}$.

The corresponding correction is
\[
    d^{(k)}
    =
    \widetilde{x}^{(k+1)}-x^{(k)}.
\]

Rather than accepting the entire
correction, we introduce a scalar
parameter $\omega$ and define
\begin{equation}
    x^{(k+1)}
    =
    x^{(k)}
    +
    \omega d^{(k)}.
    \label{eq:relaxation_correction}
\end{equation}

Substituting the definition of
$d^{(k)}$ gives
\[
    x^{(k+1)}
    =
    x^{(k)}
    +
    \omega
    \left(
        \widetilde{x}^{(k+1)}
        -
        x^{(k)}
    \right).
\]

Rearranging,
\begin{equation}
    \boxed{
    x^{(k+1)}
    =
    (1-\omega)x^{(k)}
    +
    \omega\widetilde{x}^{(k+1)}.
    }
    \label{eq:relaxation_general}
\end{equation}

The scalar $\omega$ is called the
\emph{relaxation parameter}.

Three cases are particularly important.

\begin{enumerate}
    \item \textbf{Under-relaxation:
    $0<\omega<1$.}

    We accept only a fraction of the
    suggested correction. The new
    approximation lies between the
    previous approximation and the
    candidate approximation.

    \item \textbf{No relaxation:
    $\omega=1$.}

    We accept the candidate approximation
    without modification:
    \[
        x^{(k+1)}
        =
        \widetilde{x}^{(k+1)}.
    \]

    \item \textbf{Over-relaxation:
    $\omega>1$.}

    We move beyond the candidate
    approximation in the direction
    of the suggested correction.
    This may accelerate convergence,
    although an excessively large
    parameter may cause divergence.
\end{enumerate}

For example, suppose the current
position is
\[
    x^{(k)}
    =
    \begin{bmatrix}
        0\\0
    \end{bmatrix},
\]
and the candidate position is
\[
    \widetilde{x}^{(k+1)}
    =
    \begin{bmatrix}
        1\\1
    \end{bmatrix}.
\]

Then the relaxed update becomes
\[
    x^{(k+1)}
    =
    \omega
    \begin{bmatrix}
        1\\1
    \end{bmatrix}.
\]

Thus,
\[
\begin{array}{c|c}
    \omega & x^{(k+1)}\\
    \hline
    0.5 & (0.5,0.5)\\
    0.8 & (0.8,0.8)\\
    1.0 & (1.0,1.0)\\
    1.5 & (1.5,1.5)
\end{array}
\]

Geometrically, the relaxation parameter
controls the distance traveled along
the proposed correction direction.

However, our ultimate objective is not
simply to travel farther or shorter
distances. It is to choose $\omega$
so that the resulting sequence converges
to the solution as efficiently as possible.


\subsection{Successive Over-Relaxation as a Modification
of Gauss--Seidel}

We now apply the relaxation idea to the
Gauss--Seidel method.

Recall that for a system
\[
    Ax=b,
\]
the Gauss--Seidel update for the
$i$th component is
\begin{equation}
    \widetilde{x}_i^{(k+1)}
    =
    \frac{1}{a_{ii}}
    \left(
        b_i
        -
        \sum_{j=1}^{i-1}
        a_{ij}x_j^{(k+1)}
        -
        \sum_{j=i+1}^{n}
        a_{ij}x_j^{(k)}
    \right).
    \label{eq:sor_gs_candidate}
\end{equation}

Here,
$\widetilde{x}_i^{(k+1)}$
denotes the candidate value obtained
from the Gauss--Seidel formula.

The essential feature is that the
most recently updated components
are used immediately.

Instead of accepting this candidate
value directly, we apply the
relaxation formula
\[
    x_i^{(k+1)}
    =
    (1-\omega)x_i^{(k)}
    +
    \omega\widetilde{x}_i^{(k+1)}.
\]

Substituting
\eqref{eq:sor_gs_candidate},
we obtain
\begin{equation}
\boxed{
\begin{aligned}
x_i^{(k+1)}
={}&
(1-\omega)x_i^{(k)}
\\
&+
\frac{\omega}{a_{ii}}
\left(
b_i
-
\sum_{j=1}^{i-1}
a_{ij}x_j^{(k+1)}
-
\sum_{j=i+1}^{n}
a_{ij}x_j^{(k)}
\right).
\end{aligned}
}
\label{eq:sor_component}
\end{equation}

This is the \emph{successive over-relaxation
(SOR) method}.

Notice that when $\omega=1$,
the first term disappears and
we recover the ordinary
Gauss--Seidel method.

Thus, Gauss--Seidel is a special
case of SOR.

The parameter $\omega$ provides
an additional degree of freedom
that allows us to modify the
convergence behavior without
fundamentally changing the
computational structure of
Gauss--Seidel.


\subsection{A Simple Numerical Example}

Let us return to the system
\[
\begin{aligned}
    4x+y&=9,\\
    x+3y&=7.
\end{aligned}
\]

Recall that the exact solution is
\[
    x^*=\frac{20}{11},
    \qquad
    y^*=\frac{19}{11}.
\]

The SOR iteration becomes
\begin{equation}
\begin{aligned}
x^{(k+1)}
&=
(1-\omega)x^{(k)}
+
\frac{\omega}{4}
\left(9-y^{(k)}\right),
\\[4pt]
y^{(k+1)}
&=
(1-\omega)y^{(k)}
+
\frac{\omega}{3}
\left(7-x^{(k+1)}\right).
\end{aligned}
\label{eq:sor_simple_example}
\end{equation}

Starting from
\[
    x^{(0)}=0,
    \qquad
    y^{(0)}=0,
\]
suppose we choose
\[
    \omega=1.2.
\]

The first component becomes
\[
\begin{aligned}
x^{(1)}
&=
(1-1.2)(0)
+
\frac{1.2}{4}(9-0)
\\
&=2.7.
\end{aligned}
\]

Using this newly computed value,
the second component becomes
\[
\begin{aligned}
y^{(1)}
&=
(1-1.2)(0)
+
\frac{1.2}{3}(7-2.7)
\\
&=1.72.
\end{aligned}
\]

Thus, after the first iteration,
\[
    \boxed{
    \begin{bmatrix}
        x^{(1)}\\y^{(1)}
    \end{bmatrix}
    =
    \begin{bmatrix}
        2.7\\1.72
    \end{bmatrix}.
    }
\]

Notice that the first component
has moved beyond the ordinary
Gauss--Seidel candidate $2.25$.

This is precisely the effect
of over-relaxation.

However, the value $1.2$ was
chosen only for illustration.
It is not necessarily the
optimal relaxation parameter
for this system.

Determining an optimal parameter
requires us to understand how
$\omega$ influences the
iteration matrix and its
eigenvalues.


\subsection{Matrix Formulation of SOR}

We now derive the matrix form
of the SOR iteration.

As before, write
\[
    A=D+L+U,
\]
where $D$, $L$, and $U$
denote the diagonal, strictly
lower triangular, and strictly
upper triangular parts of $A$,
respectively.

Starting from the componentwise
SOR formula
\eqref{eq:sor_component},
multiply both sides by $a_{ii}$:
\[
\begin{aligned}
a_{ii}x_i^{(k+1)}
={}&
(1-\omega)a_{ii}x_i^{(k)}
+\omega b_i
\\
&-
\omega\sum_{j<i}
a_{ij}x_j^{(k+1)}
-
\omega\sum_{j>i}
a_{ij}x_j^{(k)}.
\end{aligned}
\]

Move the terms involving
the newly computed components
to the left-hand side:
\[
\begin{aligned}
a_{ii}x_i^{(k+1)}
+
\omega\sum_{j<i}
a_{ij}x_j^{(k+1)}
={}&
(1-\omega)a_{ii}x_i^{(k)}
\\
&-
\omega\sum_{j>i}
a_{ij}x_j^{(k)}
+\omega b_i.
\end{aligned}
\]

Collecting these equations
for all components gives
\begin{equation}
    (D+\omega L)x^{(k+1)}
    =
    \left[(1-\omega)D-\omega U\right]
    x^{(k)}
    +
    \omega b.
    \label{eq:sor_matrix_intermediate}
\end{equation}

Since $D+\omega L$ is lower
triangular with nonzero
diagonal entries, it is
invertible.

Therefore,
\begin{equation}
\boxed{
\begin{aligned}
x^{(k+1)}
={}&
(D+\omega L)^{-1}
\left[(1-\omega)D-\omega U\right]
x^{(k)}
\\
&+
\omega(D+\omega L)^{-1}b.
\end{aligned}
}
\label{eq:sor_matrix}
\end{equation}

Define the SOR iteration matrix
\begin{equation}
\boxed{
T_{\omega}
=
(D+\omega L)^{-1}
\left[(1-\omega)D-\omega U\right].
}
\label{eq:sor_iteration_matrix}
\end{equation}

Also, define
\[
    c_{\omega}
    =
    \omega(D+\omega L)^{-1}b.
\]

The SOR method now takes
the familiar stationary
iteration form
\begin{equation}
    \boxed{
    x^{(k+1)}
    =
    T_{\omega}x^{(k)}
    +
    c_{\omega}.
    }
    \label{eq:sor_compact}
\end{equation}

At $\omega=1$,
\[
\begin{aligned}
T_1
&=
(D+L)^{-1}
\left[(1-1)D-U\right]
\\
&=
-(D+L)^{-1}U.
\end{aligned}
\]

Thus,
\[
    \boxed{
    T_1=T_{GS},
    }
\]
confirming once again that
Gauss--Seidel is the special
case of SOR corresponding
to $\omega=1$.

We have now introduced a
family of iteration matrices
indexed by the relaxation
parameter $\omega$.

This leads to the central
mathematical question of
our investigation:

\begin{center}
\emph{
For which values of $\omega$
does the iteration converge,
and which value makes
convergence fastest?
}
\end{center}

Since the error satisfies
\[
    e^{(k+1)}
    =
    T_{\omega}e^{(k)},
\]
we obtain
\[
    e^{(k)}
    =
    T_{\omega}^{k}e^{(0)}.
\]

Consequently, the SOR method
converges for every initial
approximation precisely when
\[
    \boxed{
    \rho(T_{\omega})<1.
    }
\]

Moreover, the spectral radius
governs the asymptotic
convergence factor.

Our objective is therefore
to determine, whenever possible,
a relaxation parameter satisfying
\[
    \boxed{
    \omega_{\mathrm{opt}}
    \in
    \operatorname*{arg\,min}_{\omega>0}
    \rho(T_{\omega}).
    }
\]

The existence, uniqueness,
and characterization of such
a parameter depend on the
properties of the coefficient
matrix.

In the following section,
we investigate these questions
and derive an explicit optimal
relaxation parameter for an
important class of matrices
arising in numerical solutions
of differential equations.



\section{Convergence Analysis and Optimal Relaxation Parameters}

We have seen that the successive over-relaxation method
introduces a parameter $\omega$ into the Gauss--Seidel
iteration, allowing us to control the magnitude of each
successive correction.

However, introducing an additional parameter does not
automatically improve an iterative method. A poor choice
may slow convergence considerably or even cause an
otherwise convergent iteration to diverge.

This raises two fundamental questions:

\begin{enumerate}
    \item For which values of $\omega$ does the SOR
    method converge?

    \item Among the admissible values, which choice
    of $\omega$ produces the fastest convergence?
\end{enumerate}

The first question concerns the stability and convergence
of the iteration. The second concerns its efficiency.

Recall that the SOR iteration matrix is
\begin{equation}
    T_\omega
    =
    (D+\omega L)^{-1}
    \left[(1-\omega)D-\omega U\right].
    \label{eq:sor_matrix_convergence}
\end{equation}

The error after $k$ iterations satisfies
\[
    e^{(k)}=T_\omega^k e^{(0)}.
\]

Consequently, the iteration converges for every initial
approximation if and only if
\begin{equation}
    \boxed{\rho(T_\omega)<1.}
    \label{eq:sor_spectral_condition}
\end{equation}

Our objective is therefore to understand how the
spectral radius of $T_\omega$ depends on $\omega$.


\subsection{A Necessary Condition for SOR Convergence}

Before examining any particular coefficient matrix,
we can establish a necessary condition that applies
to every SOR iteration with a well-defined splitting.

\begin{theorem}
Let $A\in\mathbb{R}^{n\times n}$ have nonzero diagonal
entries, and let
\[
    A=D+L+U.
\]
If the SOR method converges for every initial
approximation, then
\[
    \boxed{0<\omega<2.}
\]
\end{theorem}

\begin{proof}
Recall that
\[
    T_\omega
    =
    (D+\omega L)^{-1}
    \left[(1-\omega)D-\omega U\right].
\]

Taking determinants, we obtain
\[
\det(T_\omega)
=
\frac{
\det\left((1-\omega)D-\omega U\right)
}{
\det(D+\omega L)
}.
\]

Now observe that $D+\omega L$ is lower triangular,
while $(1-\omega)D-\omega U$ is upper triangular.

Since the determinant of a triangular matrix equals
the product of its diagonal entries,
\[
    \det(D+\omega L)
    =
    \prod_{i=1}^{n}a_{ii}.
\]

Similarly,
\[
\begin{aligned}
\det\left((1-\omega)D-\omega U\right)
&=
\prod_{i=1}^{n}(1-\omega)a_{ii}\\
&=
(1-\omega)^n\prod_{i=1}^{n}a_{ii}.
\end{aligned}
\]

Therefore,
\begin{equation}
    \boxed{
    \det(T_\omega)=(1-\omega)^n.
    }
    \label{eq:sor_determinant}
\end{equation}

Let $\lambda_1,\ldots,\lambda_n$ denote the
eigenvalues of $T_\omega$, counted with
algebraic multiplicity.

We know that the determinant equals the product
of the eigenvalues:
\[
    \det(T_\omega)
    =
    \prod_{i=1}^{n}\lambda_i.
\]

Consequently,
\[
    |1-\omega|^n
    =
    \prod_{i=1}^{n}|\lambda_i|.
\]

By the definition of spectral radius,
\[
    |\lambda_i|\leq\rho(T_\omega)
\]
for every $i$.

Hence,
\[
    |1-\omega|^n
    \leq
    \rho(T_\omega)^n.
\]

Taking the $n$th root gives
\begin{equation}
    \boxed{
    \rho(T_\omega)\geq|1-\omega|.
    }
    \label{eq:sor_spectral_lower_bound}
\end{equation}

For convergence, we require
\[
    \rho(T_\omega)<1.
\]

Therefore,
\[
    |1-\omega|<1,
\]
which is equivalent to
\[
    -1<1-\omega<1.
\]

Rearranging,
\[
    \boxed{0<\omega<2.}
\]

This proves the necessary condition.
\end{proof}


\begin{remark}
The condition $0<\omega<2$ is necessary but
not sufficient for convergence of SOR for an
arbitrary coefficient matrix.

In other words, choosing $\omega$ inside this
interval does not guarantee convergence unless
additional assumptions are imposed on $A$.

However, for symmetric positive definite matrices,
the situation is considerably more favorable:
SOR converges for every relaxation parameter
satisfying $0<\omega<2$.

We shall establish this result in the next subsection.
\end{remark}


\subsection{An Immediate Consequence for Optimal Relaxation}

The preceding theorem gives us more than a restriction
on the admissible relaxation parameter.

It also establishes the lower bound
\[
    \boxed{
    \rho(T_\omega)\geq|1-\omega|.
    }
\]

This observation is important because our objective
is to minimize the spectral radius.

In particular, for $\omega>1$,
\[
    \rho(T_\omega)\geq\omega-1.
\]

Thus, increasing $\omega$ cannot improve convergence
indefinitely. As $\omega$ approaches $2$, this lower
bound approaches $1$.

For example,
\[
\begin{array}{c|c}
\omega & \text{Lower bound for }\rho(T_\omega)\\
\hline
1.2 & 0.2\\
1.5 & 0.5\\
1.8 & 0.8\\
1.9 & 0.9\\
1.99 & 0.99
\end{array}
\]

These are lower bounds, not necessarily the actual
spectral radii.

Nevertheless, they already reveal something
interesting: although over-relaxation may accelerate
Gauss--Seidel, choosing an excessively large
relaxation parameter eventually becomes
counterproductive.

The optimal parameter must balance the acceleration
obtained through extrapolation against the possibility
of introducing slowly decaying or oscillatory errors.

This leads us naturally to the next question:

\emph{Under what conditions can we guarantee
convergence throughout the interval $0<\omega<2$?}


\subsection{Convergence of SOR for Symmetric Positive Definite Matrices}

We have established that $0<\omega<2$ is a necessary
condition for convergence of the SOR method. However,
for a general matrix, this condition alone does not
guarantee convergence.

Fortunately, for symmetric positive definite matrices,
the situation is considerably more favorable. In fact,
the same condition becomes both necessary and sufficient.

Before proving this result, let us return to the
quadratic energy interpretation introduced earlier.

Suppose that $A$ is symmetric positive definite.
Then the linear system
\[
    Ax=b
\]
is equivalent to minimizing
\begin{equation}
    E(x)=\frac12 x^T A x-b^T x.
    \label{eq:sor_energy}
\end{equation}

Since
\[
    \nabla E(x)=Ax-b
\]
and
\[
    \nabla^2 E(x)=A,
\]
the functional $E$ has a unique global minimizer
$x^*=A^{-1}b$.

Now consider what happens when we change only one
component of the current approximation.

Let $e_i$ denote the $i$th standard basis vector
and consider the point
\[
    x+\alpha e_i,
\]
where $\alpha$ represents a correction to the
$i$th component.

Substituting into the energy functional gives
\[
\begin{aligned}
E(x+\alpha e_i)
&=
\frac12(x+\alpha e_i)^T
A(x+\alpha e_i)
-b^T(x+\alpha e_i).
\end{aligned}
\]

Expanding and using the symmetry of $A$,
we obtain
\[
\begin{aligned}
E(x+\alpha e_i)
={}&
\frac12x^TAx-b^Tx
\\
&+\alpha e_i^T(Ax-b)
+\frac12\alpha^2e_i^TAe_i.
\end{aligned}
\]

Therefore,
\begin{equation}
    \boxed{
    E(x+\alpha e_i)-E(x)
    =
    \alpha(Ax-b)_i
    +\frac12a_{ii}\alpha^2.
    }
    \label{eq:sor_energy_change}
\end{equation}

Since $A$ is positive definite,
\[
    a_{ii}>0.
\]

Thus, the right-hand side is a strictly convex
quadratic function of $\alpha$.

Differentiating with respect to $\alpha$,
\[
    \frac{d}{d\alpha}
    \left[
    E(x+\alpha e_i)-E(x)
    \right]
    =
    (Ax-b)_i+a_{ii}\alpha.
\]

Setting the derivative equal to zero gives
\[
    \boxed{
    \alpha_*=
    -\frac{(Ax-b)_i}{a_{ii}}.
    }
\]

This is precisely the correction obtained by
minimizing the energy with respect to the
$i$th coordinate while keeping the other
coordinates fixed.

When the coordinates are updated successively,
using the most recently available values,
we recover the Gauss--Seidel method.

Hence, for symmetric positive definite matrices,
Gauss--Seidel can be interpreted as a sequence
of exact coordinatewise energy minimizations.

But what happens if we introduce relaxation?

Instead of taking the correction $\alpha_*$,
SOR takes
\[
    \alpha=\omega\alpha_*.
\]

Substituting into
\eqref{eq:sor_energy_change}, we obtain
\[
\begin{aligned}
E(x+\omega\alpha_*e_i)-E(x)
&=
\omega\alpha_*(Ax-b)_i
+\frac12a_{ii}\omega^2\alpha_*^2.
\end{aligned}
\]

Since
\[
    (Ax-b)_i=-a_{ii}\alpha_*,
\]
we obtain
\[
\begin{aligned}
E(x+\omega\alpha_*e_i)-E(x)
&=
-\omega a_{ii}\alpha_*^2
+\frac12\omega^2a_{ii}\alpha_*^2
\\
&=
-\frac12\omega(2-\omega)
a_{ii}\alpha_*^2.
\end{aligned}
\]

Therefore,
\begin{equation}
    \boxed{
    E(x+\omega\alpha_*e_i)-E(x)
    =
    -\frac{\omega(2-\omega)}{2}
    a_{ii}\alpha_*^2.
    }
    \label{eq:sor_energy_decrease}
\end{equation}

This identity is particularly revealing.

Since $a_{ii}>0$, the energy decreases whenever
\[
    \omega(2-\omega)>0
\]
and the coordinate correction is nonzero.

But
\[
    \omega(2-\omega)>0
    \quad\Longleftrightarrow\quad
    0<\omega<2.
\]

Thus, for every $0<\omega<2$, each SOR
coordinate update decreases the energy unless
that coordinate already requires no correction.

Notice that even when $\omega>1$, so that
the method moves beyond the coordinatewise
minimizer, the energy still decreases as long
as $\omega<2$.

This gives a geometric explanation for the
convergence interval that we previously
obtained using determinants.


\begin{theorem}[Convergence of SOR for SPD Matrices]
Let $A\in\mathbb{R}^{n\times n}$ be symmetric
positive definite. Then the successive
over-relaxation method converges to the unique
solution of $Ax=b$, for every initial
approximation, if and only if
\[
    \boxed{0<\omega<2.}
\]
\end{theorem}

\begin{proof}
We have already proved that $0<\omega<2$
is necessary for convergence.

It remains to prove sufficiency.

Let $x^*$ denote the unique solution of
$Ax=b$, and define the error
\[
    e=x-x^*.
\]

Since $Ax^*=b$, expanding the energy around
$x^*$ gives
\[
\begin{aligned}
E(x)-E(x^*)
&=
\frac12(x-x^*)^TA(x-x^*).
\end{aligned}
\]

Therefore,
\begin{equation}
    \boxed{
    E(x)-E(x^*)
    =
    \frac12 e^TAe.
    }
    \label{eq:sor_energy_error}
\end{equation}

Because $A$ is symmetric positive definite,
the expression
\[
    \|e\|_A=\sqrt{e^TAe}
\]
defines a norm.

Now consider one complete SOR iteration,
which consists of updating all $n$ coordinates
successively.

By \eqref{eq:sor_energy_decrease}, each
coordinate update decreases the energy or
leaves it unchanged.

Moreover, if every coordinate update leaves
the energy unchanged, then every coordinate
correction must be zero.

In that case, the current approximation
satisfies every equation of $Ax=b$ and
therefore equals $x^*$.

Consequently, for any $x\neq x^*$,
one complete SOR iteration strictly
decreases the energy:
\[
    E(x^{(k+1)})<E(x^{(k)}).
\]

Using \eqref{eq:sor_energy_error}, this gives
\[
    \|e^{(k+1)}\|_A
    <
    \|e^{(k)}\|_A
\]
whenever $e^{(k)}\neq0$.

Recall that SOR is a stationary iteration,
so its errors satisfy
\[
    e^{(k+1)}=T_\omega e^{(k)}.
\]

Thus,
\[
    \|T_\omega e\|_A<\|e\|_A
    \qquad\text{for every }e\neq0.
\]

Since $T_\omega$ is linear and the unit
sphere
\[
    S=\{e\in\mathbb{R}^n:\|e\|_A=1\}
\]
is compact, the continuous function
\[
    e\longmapsto\|T_\omega e\|_A
\]
attains its maximum on $S$.

The strict inequality above implies that
this maximum is some constant $q<1$.

Hence,
\[
    \|T_\omega e\|_A
    \leq q\|e\|_A
\]
for every $e\in\mathbb{R}^n$.

Applying this inequality repeatedly,
\[
    \|e^{(k)}\|_A
    \leq q^k\|e^{(0)}\|_A.
\]

Since $0\leq q<1$,
\[
    \lim_{k\to\infty}\|e^{(k)}\|_A=0.
\]

Therefore,
\[
    \boxed{
    \lim_{k\to\infty}x^{(k)}=x^*.
    }
\]

Thus, SOR converges for every initial
approximation whenever $0<\omega<2$.
\end{proof}


\begin{remark}
The preceding proof establishes convergence,
but it does not determine the optimal
relaxation parameter.

Indeed, the energy decrease in
\eqref{eq:sor_energy_decrease} is largest
for an individual coordinate update when
$\omega=1$, corresponding to Gauss--Seidel.

However, minimizing the energy during one
coordinate update is not the same as minimizing
the number of complete iterations required
to solve the system.

By moving beyond the coordinatewise minimizer,
SOR may produce a sequence of approximations
that converges substantially faster in the
long run.

This is why over-relaxation can outperform
ordinary Gauss--Seidel even though each
individual Gauss--Seidel coordinate update
minimizes the energy along its coordinate.
\end{remark}


\section{Derivation of the Optimal SOR Parameter}

We have established that the SOR method converges
for every $0<\omega<2$ when the coefficient matrix
is symmetric positive definite.

However, convergence is only part of our objective.
We are particularly interested in choosing a
relaxation parameter that makes the iteration
converge as rapidly as possible.

Since the asymptotic convergence factor is
governed by the spectral radius, our problem is
to determine
\begin{equation}
    \boxed{
    \omega_{\mathrm{opt}}
    \in
    \operatorname*{arg\,min}_{0<\omega<2}
    \rho(T_\omega).
    }
    \label{eq:optimal_sor_problem}
\end{equation}

For a general matrix, this optimization problem
may be difficult to solve analytically.

Fortunately, for certain structured matrices,
the eigenvalues of the SOR iteration matrix
are related to those of the Jacobi iteration
matrix through a remarkable algebraic identity.

We shall derive this relationship first and
then use it to determine the optimal
relaxation parameter.


\subsection{The Relationship Between Jacobi and SOR Eigenvalues}

Consider a nonsingular tridiagonal matrix
\[
A=
\begin{bmatrix}
a_1 & c_1 & 0 & \cdots & 0\\
b_1 & a_2 & c_2 & \ddots & \vdots\\
0 & b_2 & a_3 & \ddots & 0\\
\vdots & \ddots & \ddots & \ddots & c_{n-1}\\
0 & \cdots & 0 & b_{n-1} & a_n
\end{bmatrix},
\]
where every diagonal entry $a_i$ is nonzero.

As before, write
\[
    A=D+L+U.
\]

The Jacobi iteration matrix is
\begin{equation}
    T_J=-D^{-1}(L+U),
    \label{eq:jacobi_matrix_optimal}
\end{equation}
while the SOR iteration matrix is
\begin{equation}
    T_\omega
    =
    (D+\omega L)^{-1}
    \left[(1-\omega)D-\omega U\right].
    \label{eq:sor_matrix_optimal}
\end{equation}

Our objective is to establish a relationship
between the eigenvalues of these two matrices.

Let $\lambda$ be an eigenvalue of $T_\omega$.
Then
\[
    \det(\lambda I-T_\omega)=0.
\]

Substituting the expression for $T_\omega$,
we obtain
\[
\det\left[
\lambda I-
(D+\omega L)^{-1}
\left((1-\omega)D-\omega U\right)
\right]=0.
\]

Multiplying the matrix inside the determinant
on the left by $D+\omega L$, which is invertible,
we obtain the equivalent condition
\[
\det\left[
\lambda(D+\omega L)
-
\left((1-\omega)D-\omega U\right)
\right]=0.
\]

Expanding,
\[
\det\left[
\lambda D+\lambda\omega L
-(1-\omega)D+\omega U
\right]=0.
\]

Collecting the diagonal terms gives
\begin{equation}
\boxed{
\det\left[
(\lambda+\omega-1)D
+\lambda\omega L+\omega U
\right]=0.
}
\label{eq:sor_characteristic_intermediate}
\end{equation}

At first sight, this expression does not seem
particularly helpful. The coefficients of $L$
and $U$ are different, and the Jacobi matrix
contains their sum $L+U$.

But here the tridiagonal structure of $A$
plays an important role.

Suppose $\lambda\neq0$, and choose a complex
number $s$ satisfying
\[
    s^2=\lambda.
\]

Consider the invertible diagonal matrix
\[
    S=
    \operatorname{diag}
    (1,s,s^2,\ldots,s^{n-1}).
\]


Because $L$ contains nonzero entries only immediately
below the diagonal, and $U$ contains nonzero entries
only immediately above it, we can examine the
similarity transformation entry by entry.

Recall that
\[
S=\operatorname{diag}(1,s,s^2,\ldots,s^{n-1}),
\]
and therefore
\[
S^{-1}=\operatorname{diag}
(1,s^{-1},s^{-2},\ldots,s^{-(n-1)}).
\]

Since both $S$ and $S^{-1}$ are diagonal, for any
matrix $M$ we have
\[
(S^{-1}MS)_{ij}
=
(S^{-1})_{ii}M_{ij}S_{jj}
=
s^{-(i-1)}M_{ij}s^{j-1}.
\]

Thus,
\begin{equation}
\boxed{
(S^{-1}MS)_{ij}=s^{j-i}M_{ij}.
}
\end{equation}

Now consider $L$. Its only potentially nonzero
entries occur when $i=j+1$. Hence,
\[
(S^{-1}LS)_{ij}
=
s^{j-i}L_{ij}
=
s^{-1}L_{ij}.
\]

All other entries of $L$ are zero and remain zero.
Consequently,
\[
\boxed{S^{-1}LS=s^{-1}L.}
\]

Similarly, the only potentially nonzero entries
of $U$ occur when $j=i+1$. Therefore,
\[
(S^{-1}US)_{ij}
=
s^{j-i}U_{ij}
=
sU_{ij}.
\]

Hence,
\[
\boxed{S^{-1}US=sU.}
\]

Finally, since $D$ and $S$ are diagonal matrices,
they commute, and therefore
\[
\boxed{S^{-1}DS=D.}
\]



Consequently,
\[
\begin{aligned}
&S^{-1}
\left[
(\lambda+\omega-1)D
+\lambda\omega L+\omega U
\right]S
\\[4pt]
&=
(\lambda+\omega-1)D
+\lambda\omega s^{-1}L
+\omega sU.
\end{aligned}
\]

Using $\lambda=s^2$, we obtain
\[
    \lambda s^{-1}=s.
\]

Therefore,
\[
\begin{aligned}
&S^{-1}
\left[
(\lambda+\omega-1)D
+\lambda\omega L+\omega U
\right]S
\\
&=
(\lambda+\omega-1)D
+\omega s(L+U).
\end{aligned}
\]

Since similar matrices have the same determinant,
equation
\eqref{eq:sor_characteristic_intermediate}
becomes
\begin{equation}
\boxed{
\det\left[
(\lambda+\omega-1)D
+\omega s(L+U)
\right]=0.
}
\label{eq:sor_jacobi_bridge}
\end{equation}

Now factor out $\omega sD$:
\[
\det\left[
\omega sD
\left(
\frac{\lambda+\omega-1}{\omega s}I
+
D^{-1}(L+U)
\right)
\right]=0.
\]

Since $\omega sD$ is invertible, we must have
\[
\det\left[
\frac{\lambda+\omega-1}{\omega s}I
+
D^{-1}(L+U)
\right]=0.
\]

Recall that
\[
    T_J=-D^{-1}(L+U).
\]

Hence,
\[
\det\left[
\frac{\lambda+\omega-1}{\omega s}I
-T_J
\right]=0.
\]

This means that
\[
    \mu=
    \frac{\lambda+\omega-1}{\omega s}
\]
is an eigenvalue of the Jacobi iteration matrix.

Therefore,
\[
    \lambda+\omega-1=\omega\mu s.
\]

Squaring both sides and recalling that
$s^2=\lambda$, we obtain
\begin{equation}
\boxed{
    (\lambda+\omega-1)^2
    =
    \omega^2\mu^2\lambda.
}
\label{eq:young_eigenvalue_relation}
\end{equation}

This is the fundamental relationship between
the eigenvalues of the Jacobi and SOR iteration
matrices for the tridiagonal systems under
consideration.

The derivation assumed $\lambda\neq0$.
The zero-eigenvalue case can be treated
separately: from
\[
    \det(T_\omega)=(1-\omega)^n,
\]
we see that $\lambda=0$ can occur only
when $\omega=1$, in which case the
relationship also holds with $\mu=0$.

The significance of
\eqref{eq:young_eigenvalue_relation}
is that it allows us to study the eigenvalues
of SOR through those of Jacobi.

Rather than computing the spectrum of
$T_\omega$ independently for every
relaxation parameter, we can exploit the
simpler structure of $T_J$.

In the next subsection, we shall use this
identity to determine the value of $\omega$
that minimizes the spectral radius.

\subsection{Derivation of the Optimal Relaxation Parameter}

We have established the relationship
\begin{equation}
    (\lambda+\omega-1)^2
    =
    \omega^2\mu^2\lambda,
    \label{eq:sor_jacobi_eigenvalue_relation}
\end{equation}
where $\mu$ is an eigenvalue of the Jacobi
iteration matrix and $\lambda$ is an
associated eigenvalue of the SOR iteration
matrix.

Our objective is to determine the value of
$\omega$ that minimizes the spectral radius
of $T_\omega$.

For the remainder of this derivation, we
assume that $A$ is symmetric positive definite
and tridiagonal, with nonzero diagonal entries,
and that the Jacobi iteration converges.

Consequently, the eigenvalues of $T_J$ are
real, and
\[
    r:=\rho(T_J)<1.
\]

We shall see that the optimal relaxation
parameter depends only on $r$.

\subsubsection*{A Quadratic Equation for the SOR Eigenvalues}

Expanding \eqref{eq:sor_jacobi_eigenvalue_relation},
we obtain
\[
    \lambda^2
    +2(\omega-1)\lambda
    +(\omega-1)^2
    =
    \omega^2\mu^2\lambda.
\]

Collecting terms,
\begin{equation}
\boxed{
    \lambda^2
    +
    \left[
        2(\omega-1)-\omega^2\mu^2
    \right]\lambda
    +
    (\omega-1)^2
    =0.
}
\label{eq:sor_eigenvalue_quadratic}
\end{equation}

Thus, the eigenvalues of SOR are determined
by a quadratic equation involving the
relaxation parameter $\omega$ and the
Jacobi eigenvalues $\mu$.

Recall from elementary algebra that the
nature of the roots of a quadratic equation
is determined by its discriminant.

For \eqref{eq:sor_eigenvalue_quadratic},
the discriminant is
\[
\begin{aligned}
\Delta
&=
\left[
2(\omega-1)-\omega^2\mu^2
\right]^2
-4(\omega-1)^2
\\
&=
\omega^4\mu^4
-4\omega^2\mu^2(\omega-1).
\end{aligned}
\]

Factoring,
\begin{equation}
\boxed{
    \Delta
    =
    \omega^2\mu^2
    \left[
        \omega^2\mu^2-4(\omega-1)
    \right].
}
\label{eq:sor_discriminant}
\end{equation}

This expression is the key to determining
the optimal relaxation parameter.


\subsubsection*{What Happens When the Discriminant Changes Sign?}

Consider first the case $\Delta>0$.

The quadratic equation has two distinct
real roots. Their magnitudes depend on
$\omega$ and $\mu$, and one of them may
dominate the convergence behavior.
%```latex
To see why,
consider a Jacobi eigenvalue satisfying
$|\mu|=r=\rho(T_J)$. For
$0<\omega<\omega_*$, where $\omega_*$ denotes
the transition parameter to be determined
shortly, the corresponding discriminant
is strictly positive. Solving the quadratic
equation gives two distinct real SOR
eigenvalues:
\[
\lambda_{\pm}
=
\frac{
\omega^2r^2-2(\omega-1)
\pm
\omega r
\sqrt{\omega^2r^2-4(\omega-1)}
}{2}.
\]
The larger eigenvalue $\lambda_+$ is
positive and provides a lower bound for
the spectral radius:
\[
    \rho(T_\omega)\geq\lambda_+.
\]
As we shall prove later, this eigenvalue
decreases as $\omega$ approaches
$\omega_*$ from below. At the transition,
the two real eigenvalues coincide, and
their common value is $\omega_*-1$.
Beyond this point, the eigenvalues
associated with $|\mu|=r$ become complex
conjugates, and their common magnitude
increases with $\omega$. Thus, the
positive-discriminant case explains
why increasing $\omega$ initially
improves convergence, while the
negative-discriminant case explains
why increasing it beyond the transition
eventually becomes counterproductive.
Together, these two behaviors lead
us to the optimal relaxation parameter.


Now suppose $\Delta<0$.

The roots become a complex conjugate pair,
say $\lambda_1$ and $\lambda_2$.

By Vieta's formula,
\[
    \lambda_1\lambda_2=(\omega-1)^2.
\]

Since the roots are complex conjugates,
\[
    \lambda_2=\overline{\lambda_1}.
\]

Therefore,
\[
    |\lambda_1|^2=(\omega-1)^2,
\]
and hence
\begin{equation}
    \boxed{
    |\lambda_1|=|\lambda_2|=|\omega-1|.
    }
\end{equation}

This is particularly interesting.

Recall that we previously established
the universal lower bound
\[
    \rho(T_\omega)\geq|\omega-1|.
\]

Thus, when the relevant SOR eigenvalues
form complex conjugate pairs, their
magnitudes attain this lower bound.

However, we must consider all Jacobi
eigenvalues, not merely one of them.

The largest magnitude of a Jacobi
eigenvalue is
\[
    r=\rho(T_J).
\]

For $\omega>1$, the expression
\[
    \omega^2\mu^2-4(\omega-1)
\]
is largest when $|\mu|=r$.

Consequently, the transition at which
the discriminant becomes nonpositive
for every Jacobi eigenvalue is determined by
\begin{equation}
    \boxed{
    \omega^2r^2-4(\omega-1)=0.
    }
    \label{eq:sor_optimal_condition}
\end{equation}

At this transition, the two roots
associated with $|\mu|=r$ coincide.

We shall now solve this equation.


\subsubsection*{Solving for the Optimal Parameter}

Rearranging
\eqref{eq:sor_optimal_condition},
we obtain
\[
    r^2\omega^2-4\omega+4=0.
\]

Applying the quadratic formula,
\[
    \omega
    =
    \frac{
        4\pm\sqrt{16-16r^2}
    }{2r^2}.
\]

Simplifying,
\[
    \omega
    =
    \frac{
        2\left(1\pm\sqrt{1-r^2}\right)
    }{r^2}.
\]

Since
\[
    r^2=
    \left(1-\sqrt{1-r^2}\right)
    \left(1+\sqrt{1-r^2}\right),
\]
the two solutions can be written as
\[
    \omega_-
    =
    \frac{2}{1+\sqrt{1-r^2}},
\]
and
\[
    \omega_+
    =
    \frac{2}{1-\sqrt{1-r^2}}.
\]

For $0<r<1$, the second solution
satisfies $\omega_+>2$.

But we already know that SOR can converge
only when
\[
    0<\omega<2.
\]

Therefore, the admissible transition
parameter is
\begin{equation}
\boxed{
    \omega_*
    =
    \frac{2}{1+\sqrt{1-r^2}}.
}
\label{eq:sor_transition_parameter}
\end{equation}

Notice that
\[
    1<\omega_*<2.
\]

Thus, the transition occurs in the
over-relaxation region.

We still need to verify that this
transition actually minimizes the
spectral radius.




\subsubsection*{Why Is This Parameter Optimal?}

We have obtained the candidate
\[
    \omega_*=
    \frac{2}{1+\sqrt{1-r^2}},
    \qquad r=\rho(T_J), \quad 0<r<1.
\]

However, finding a parameter by setting the
discriminant equal to zero does not automatically
prove that the parameter is optimal.

To establish optimality, we must show that
\[
    \rho(T_\omega)>\rho(T_{\omega_*})
\]
for every $\omega\neq\omega_*$ in the
convergence interval $0<\omega<2$.

We shall examine two cases.

\medskip
\noindent
\textbf{Case 1: $\omega_*\leq\omega<2$.}

Recall that the discriminant of the quadratic
equation for the SOR eigenvalues is
\[
    \Delta=
    \omega^2\mu^2
    \left[\omega^2\mu^2-4(\omega-1)\right].
\]

Since every Jacobi eigenvalue satisfies
$|\mu|\leq r$, we have
\[
    \omega^2\mu^2-4(\omega-1)
    \leq
    \omega^2r^2-4(\omega-1).
\]

The right-hand side vanishes at
$\omega=\omega_*$ and is negative for
$\omega_*<\omega<2$. Therefore,
\[
    \boxed{\Delta\leq0}
\]
for every Jacobi eigenvalue throughout
this interval.

When $\Delta<0$, the two roots of the
quadratic equation form a complex conjugate
pair, say $\lambda_1$ and $\lambda_2$.

By Vieta's formula,
\[
    \lambda_1\lambda_2=(\omega-1)^2.
\]

Since $\lambda_2=\overline{\lambda_1}$,
\[
    |\lambda_1|^2=(\omega-1)^2.
\]

Thus,
\[
    |\lambda_1|=|\lambda_2|=\omega-1.
\]

The same conclusion holds when $\Delta=0$,
because the repeated root is then
$\lambda=\omega-1$.

Consequently, every SOR eigenvalue has
magnitude $\omega-1$, and hence
\begin{equation}
    \boxed{
    \rho(T_\omega)=\omega-1,
    \qquad \omega_*\leq\omega<2.
    }
    \label{eq:sor_rho_right}
\end{equation}

But $\omega-1$ is strictly increasing
as a function of $\omega$.

Therefore,
\[
    \rho(T_\omega)>
    \rho(T_{\omega_*})
    \qquad
    \text{whenever }\omega>\omega_*.
\]

We have proved that choosing a parameter
larger than $\omega_*$ cannot improve
the asymptotic convergence factor.

\medskip
\noindent
\textbf{Case 2: $0<\omega<\omega_*$.}

Choose a Jacobi eigenvalue satisfying
\[
    |\mu|=r.
\]

For this eigenvalue, the relationship
between Jacobi and SOR becomes
\[
    (\lambda+\omega-1)^2
    =
    \omega^2r^2\lambda.
\]

We may write
\[
    \lambda+\omega-1
    =
    \omega r\sqrt{\lambda}
\]
for the branch corresponding to the
larger positive root.

Introducing
\[
    t=\sqrt{\lambda},
\]
we obtain
\[
    t^2+\omega-1=\omega rt,
\]
or equivalently,
\begin{equation}
    \boxed{
    t^2-\omega rt+\omega-1=0.
    }
    \label{eq:sor_t_quadratic}
\end{equation}

Applying the quadratic formula gives
\[
    t_{\pm}(\omega)
    =
    \frac{
        \omega r
        \pm
        \sqrt{\omega^2r^2-4(\omega-1)}
    }{2}.
\]

We shall concentrate on the larger root,
\begin{equation}
    \boxed{
    t_+(\omega)
    =
    \frac{
        \omega r+
        \sqrt{\omega^2r^2-4(\omega-1)}
    }{2}.
    }
    \label{eq:sor_larger_root}
\end{equation}

For $0<\omega<\omega_*$, the expression
under the square root is strictly positive.
Thus, $t_+(\omega)$ is real and positive.

Since $\lambda=t_+(\omega)^2$ is an
eigenvalue of the SOR iteration matrix,
we have
\begin{equation}
    \boxed{
    \rho(T_\omega)\geq t_+(\omega)^2.
    }
    \label{eq:sor_rho_lower_root}
\end{equation}

Therefore, it is sufficient to prove that
$t_+(\omega)$ decreases as $\omega$
approaches $\omega_*$.

\medskip
\noindent
\textbf{Showing that $t_+(\omega)$ decreases.}

Rather than differentiating the explicit
square-root expression, we can differentiate
the quadratic equation
\eqref{eq:sor_t_quadratic} implicitly.

Recall that
\[
    t_+^2-\omega rt_++\omega-1=0,
\]
where $t_+$ is a function of $\omega$.

Differentiating both sides with respect
to $\omega$,
\[
    2t_+\frac{dt_+}{d\omega}
    -rt_+
    -\omega r\frac{dt_+}{d\omega}
    +1=0.
\]

Collecting the derivative terms,
\[
    (2t_+-\omega r)
    \frac{dt_+}{d\omega}
    =
    rt_+-1.
\]

Hence,
\begin{equation}
    \boxed{
    \frac{dt_+}{d\omega}
    =
    \frac{rt_+-1}{2t_+-\omega r}.
    }
    \label{eq:sor_root_derivative}
\end{equation}

We now examine the signs of the numerator
and denominator separately.

First, from
\eqref{eq:sor_larger_root},
\[
    2t_+-\omega r
    =
    \sqrt{\omega^2r^2-4(\omega-1)}.
\]

Since $0<\omega<\omega_*$, the
discriminant is strictly positive.
Therefore,
\[
    \boxed{2t_+-\omega r>0.}
\]

Next, we must show that
\[
    rt_+-1<0.
\]

To establish this, consider the quadratic
polynomial
\[
    p(t)=t^2-\omega rt+\omega-1.
\]

At $t=1$,
\[
\begin{aligned}
    p(1)
    &=1-\omega r+\omega-1\\
    &=\omega(1-r)>0,
\end{aligned}
\]
because $0<r<1$.

Moreover, the vertex of this quadratic
occurs at
\[
    t=\frac{\omega r}{2}<1,
\]
since $0<\omega<\omega_*<2$.

The quadratic opens upward and has two
real roots. Since its vertex lies to the
left of $1$ and $p(1)>0$, both roots
must lie to the left of $1$.

In particular,
\[
    \boxed{0<t_+<1.}
\]

Since $0<r<1$, it follows that
\[
    rt_+<1,
\]
and therefore
\[
    \boxed{rt_+-1<0.}
\]

Returning to
\eqref{eq:sor_root_derivative},
we have
\[
    \frac{dt_+}{d\omega}
    =
    \frac{
        \underbrace{rt_+-1}_{<0}
    }{
        \underbrace{2t_+-\omega r}_{>0}
    }.
\]

Consequently,
\begin{equation}
    \boxed{
    \frac{dt_+}{d\omega}<0,
    \qquad 0<\omega<\omega_*.
    }
    \label{eq:sor_root_decreasing}
\end{equation}

Thus, $t_+(\omega)$ is strictly decreasing
throughout the interval $(0,\omega_*)$.

\medskip
\noindent
\textbf{Comparing the roots at the transition.}

At $\omega=\omega_*$, the discriminant
vanishes:
\[
    \omega_*^2r^2-4(\omega_*-1)=0.
\]

Therefore,
\[
    \omega_*^2r^2=4(\omega_*-1).
\]

The two roots coincide at
\[
    t_+(\omega_*)=\frac{\omega_*r}{2}.
\]

Squaring,
\[
\begin{aligned}
    t_+(\omega_*)^2
    &=
    \frac{\omega_*^2r^2}{4}\\
    &=\omega_*-1.
\end{aligned}
\]

Hence,
\begin{equation}
    \boxed{
    t_+(\omega_*)^2=\omega_*-1.
    }
    \label{eq:sor_root_transition}
\end{equation}

Since $t_+(\omega)$ is strictly
decreasing for $0<\omega<\omega_*$,
we have
\[
    t_+(\omega)>t_+(\omega_*).
\]

Both quantities are positive, so
\[
    t_+(\omega)^2>
    t_+(\omega_*)^2
    =
    \omega_*-1.
\]

Combining this inequality with
\eqref{eq:sor_rho_lower_root},
we conclude that
\begin{equation}
    \boxed{
    \rho(T_\omega)>
    \omega_*-1,
    \qquad 0<\omega<\omega_*.
    }
    \label{eq:sor_rho_left}
\end{equation}

On the other hand, from
\eqref{eq:sor_rho_right},
\[
    \rho(T_{\omega_*})=\omega_*-1.
\]

Therefore,
\[
    \rho(T_\omega)>
    \rho(T_{\omega_*})
    \qquad
    \text{for every }0<\omega<\omega_*.
\]

This completes the second case.

\medskip
\noindent
\textbf{Conclusion.}

We have proved that
\[
    \rho(T_\omega)>
    \rho(T_{\omega_*})
\]
whenever
\[
    0<\omega<2,
    \qquad \omega\neq\omega_*.
\]

Thus, the spectral radius attains its
minimum at $\omega=\omega_*$.

Recalling that
\[
    r=\rho(T_J),
\]
we finally obtain
\begin{equation}
    \boxed{
    \omega_{\mathrm{opt}}
    =
    \frac{2}{
        1+\sqrt{1-\rho(T_J)^2}
    }.
    }
    \label{eq:sor_optimal_omega}
\end{equation}

Furthermore, the corresponding optimal
spectral radius is
\begin{equation}
    \boxed{
    \rho(T_{\omega_{\mathrm{opt}}})
    =
    \omega_{\mathrm{opt}}-1.
    }
    \label{eq:sor_optimal_spectral_radius}
\end{equation}

This result provides an explicit connection
between the convergence properties of the
Jacobi method and the optimal relaxation
parameter of SOR.

It is worth emphasizing that this formula
does not hold for every symmetric positive
definite matrix. Our derivation relies on
the special eigenvalue relationship
established for the tridiagonal systems
under consideration.

For more general matrices, additional
structural assumptions are required to
obtain the same formula.


\subsection{Optimal SOR Parameter for a Finite-Difference Matrix}

We have established that, for the symmetric positive
definite tridiagonal matrices under consideration,
the optimal SOR relaxation parameter is
\[
    \omega_{\mathrm{opt}}
    =
    \frac{2}{1+\sqrt{1-\rho(T_J)^2}},
\]
where $T_J$ denotes the Jacobi iteration matrix.

The formula is theoretically useful, but it still
requires knowledge of the spectral radius of $T_J$.
For a general matrix, determining this quantity
may itself be computationally expensive.

Fortunately, the matrices arising from the
finite-difference discretization of certain
boundary value problems possess a special
structure that allows us to calculate their
eigenvalues explicitly.

\subsubsection*{Construction of the Finite-Difference Matrix}

Consider the boundary value problem
\[
    -u''(x)=f(x),
    \qquad 0<x<1,
\]
subject to the homogeneous Dirichlet
boundary conditions
\[
    u(0)=u(1)=0.
\]

Divide $[0,1]$ into $n+1$ equal subintervals,
with grid spacing
\[
    h=\frac{1}{n+1}.
\]

The interior grid points are
\[
    x_i=ih,
    \qquad i=1,\ldots,n.
\]

Using the central-difference approximation,
\[
    u''(x_i)
    \approx
    \frac{u_{i-1}-2u_i+u_{i+1}}{h^2},
\]
we obtain
\[
    -u_{i-1}+2u_i-u_{i+1}
    =
    h^2f(x_i).
\]

Thus, the finite-difference discretization
produces the linear system
\[
    Au=b,
\]
where
\[
    A=
    \begin{pmatrix}
    2 & -1 & 0 & \cdots & 0\\
    -1 & 2 & -1 & \ddots & \vdots\\
    0 & -1 & 2 & \ddots & 0\\
    \vdots & \ddots & \ddots & \ddots & -1\\
    0 & \cdots & 0 & -1 & 2
    \end{pmatrix}_{n\times n},
\]
and
\[
    b_i=h^2f(x_i).
\]

The matrix $A$ is symmetric positive definite
and tridiagonal. Hence, the optimal SOR
parameter derived previously applies.


\subsubsection*{Eigenvalues of the Jacobi Iteration Matrix}

For the finite-difference matrix,
\[
    D=2I.
\]

Recall that the Jacobi iteration matrix is
\[
    T_J=-D^{-1}(L+U).
\]

Therefore,
\[
    T_J=
    \frac12
    \begin{pmatrix}
    0 & 1 & 0 & \cdots & 0\\
    1 & 0 & 1 & \ddots & \vdots\\
    0 & 1 & 0 & \ddots & 0\\
    \vdots & \ddots & \ddots & \ddots & 1\\
    0 & \cdots & 0 & 1 & 0
    \end{pmatrix}.
\]

We now determine the eigenvalues of $T_J$
analytically.

Let
\[
    T_Jv=\mu v,
\]
where $v\neq0$ is an eigenvector.

Writing this equation componentwise gives
\[
    \frac12(v_{i-1}+v_{i+1})
    =
    \mu v_i.
\]

Introducing the boundary values
\[
    v_0=v_{n+1}=0,
\]
we can express all $n$ equations through
the recurrence relation
\begin{equation}
    v_{i-1}+v_{i+1}
    =
    2\mu v_i,
    \qquad i=1,\ldots,n.
    \label{eq:jacobi_eigenvector_recurrence}
\end{equation}

This is a second-order linear difference
equation, similar in structure to the
differential equation from which the
finite-difference matrix originated.

\medskip
\noindent
\textbf{Solving the characteristic equation.}

As with linear differential equations
with constant coefficients, we look for
solutions of the form
\[
    v_i=z^i,
\]
where $z$ is a nonzero constant to be
determined.

Substituting into
\eqref{eq:jacobi_eigenvector_recurrence},
we obtain
\[
    z^{i-1}+z^{i+1}=2\mu z^i.
\]

Dividing by $z^{i-1}$ gives
\[
    1+z^2=2\mu z.
\]

Rearranging,
\begin{equation}
    \boxed{
    z^2-2\mu z+1=0.
    }
    \label{eq:jacobi_characteristic_equation}
\end{equation}

Applying the quadratic formula,
\[
    z_{1,2}
    =
    \frac{
        2\mu\pm\sqrt{4\mu^2-4}
    }{2}.
\]

Factoring $4$ inside the square root,
\[
    z_{1,2}
    =
    \frac{
        2\mu\pm2\sqrt{\mu^2-1}
    }{2}.
\]

Therefore,
\begin{equation}
    \boxed{
    z_{1,2}
    =
    \mu\pm\sqrt{\mu^2-1}.
    }
    \label{eq:jacobi_characteristic_roots}
\end{equation}

Before proceeding, we should determine
whether these roots are real or complex.

Since $T_J$ is real symmetric, its
eigenvalues $\mu$ are real.

Furthermore, we claim that
\[
    -1<\mu<1.
\]

Indeed, the maximum absolute row-sum
of $T_J$ is at most $1$, so
\[
    |\mu|\leq\rho(T_J)\leq1.
\]

Equality cannot occur. To see this,
suppose $|\mu|=1$ and choose an
eigenvector component $v_i$ of maximum
absolute value. The recurrence
\[
    2\mu v_i=v_{i-1}+v_{i+1}
\]
shows that equality in the triangle
inequality forces neighboring
components to have the same maximum
absolute value. Repeating this
argument toward an endpoint eventually
contradicts the boundary condition
$v_0=v_{n+1}=0$.

Thus,
\[
    \boxed{-1<\mu<1.}
\]

Consequently,
\[
    \mu^2-1<0.
\]

The characteristic roots are therefore
complex conjugates. We can rewrite them as
\[
    z_{1,2}
    =
    \mu\pm i\sqrt{1-\mu^2},
\]
where $i^2=-1$.

Notice something interesting:
\[
\begin{aligned}
    |z_{1,2}|^2
    &=
    \mu^2+(1-\mu^2)\\
    &=1.
\end{aligned}
\]

Hence,
\[
    \boxed{|z_1|=|z_2|=1.}
\]

Thus, both characteristic roots lie on
the unit circle in the complex plane.

\medskip
\noindent
\textbf{Introducing the trigonometric parameter.}

Because $-1<\mu<1$, there exists
a unique angle $\theta\in(0,\pi)$
such that
\[
    \boxed{\mu=\cos\theta.}
\]

This substitution is useful because
\[
\begin{aligned}
    \sqrt{1-\mu^2}
    &=
    \sqrt{1-\cos^2\theta}\\
    &=
    \sqrt{\sin^2\theta}\\
    &=\sin\theta.
\end{aligned}
\]

Here we used the fact that
$\sin\theta>0$ for $0<\theta<\pi$.

Substituting into the characteristic
roots gives
\[
    z_{1,2}
    =
    \cos\theta\pm i\sin\theta.
\]

Recall Euler's formula,
\[
    e^{i\theta}
    =
    \cos\theta+i\sin\theta,
\]
and similarly,
\[
    e^{-i\theta}
    =
    \cos\theta-i\sin\theta.
\]

Therefore,
\begin{equation}
    \boxed{
    z_1=e^{i\theta},
    \qquad
    z_2=e^{-i\theta}.
    }
    \label{eq:jacobi_complex_roots}
\end{equation}

We now see why introducing
$\mu=\cos\theta$ was natural:
the characteristic roots lie on
the unit circle, and the angle
$\theta$ provides a convenient
way to represent them.

\medskip
\noindent
\textbf{Constructing the general solution.}

Since the characteristic equation has
two distinct roots, the general
solution of the second-order recurrence
is
\[
    v_k=C_1z_1^k+C_2z_2^k.
\]

Substituting
\eqref{eq:jacobi_complex_roots},
we obtain
\[
    v_k=
    C_1e^{ik\theta}
    +
    C_2e^{-ik\theta}.
\]



Using Euler's formula again,
\[
\begin{aligned}
    v_k
    &=
    C_1\bigl(
        \cos(k\theta)
        +\mathrm{i}\sin(k\theta)
    \bigr)\\
    &\quad+
    C_2\bigl(
        \cos(k\theta)
        -\mathrm{i}\sin(k\theta)
    \bigr).
\end{aligned}
\]

Collecting the cosine and sine terms,
\[
    v_k=
    (C_1+C_2)\cos(k\theta)
    +
    \mathrm{i}(C_1-C_2)\sin(k\theta).
\]

Writing
\[
    C=C_1+C_2,
    \qquad
    S=\mathrm{i}(C_1-C_2),
\]
we obtain the equivalent real form
\begin{equation}
    \boxed{
    v_k=C\cos(k\theta)+S\sin(k\theta).
    }
    \label{eq:jacobi_general_eigenvector}
\end{equation}

The constants $C$ and $S$ may be
chosen real because $T_J$ is a
real symmetric matrix.

\medskip
\noindent
\textbf{Applying the boundary conditions.}

Recall that our eigenvector must satisfy
\[
    v_0=v_{n+1}=0.
\]

Applying the first boundary condition
to \eqref{eq:jacobi_general_eigenvector},
we obtain
\[
\begin{aligned}
    v_0
    &=C\cos(0)+S\sin(0)\\
    &=C.
\end{aligned}
\]

Since $v_0=0$, we must have
\[
    C=0.
\]

Therefore, the eigenvector components
reduce to
\[
    v_k=S\sin(k\theta).
\]

Now apply the second boundary condition:
\[
    v_{n+1}
    =
    S\sin((n+1)\theta)=0.
\]

Since the eigenvector is nonzero,
we must have $S\neq0$.

Consequently,
\[
    \sin((n+1)\theta)=0.
\]

Recall that
\[
    \sin\alpha=0
    \quad\Longleftrightarrow\quad
    \alpha=j\pi,
    \qquad j\in\mathbb{Z}.
\]

Therefore,
\[
    (n+1)\theta=j\pi,
\]
which gives
\[
    \theta=\frac{j\pi}{n+1}.
\]

Because $0<\theta<\pi$, the admissible
integers are
\[
    j=1,2,\ldots,n.
\]

Thus, the $n$ admissible angles are
\begin{equation}
    \boxed{
    \theta_j=\frac{j\pi}{n+1},
    \qquad j=1,\ldots,n.
    }
\end{equation}

Finally, recalling that
$\mu=\cos\theta$, we obtain the
eigenvalues of the Jacobi iteration
matrix:
\begin{equation}
    \boxed{
    \mu_j=
    \cos\left(\frac{j\pi}{n+1}\right),
    \qquad j=1,\ldots,n.
    }
    \label{eq:jacobi_fd_eigenvalues}
\end{equation}

The corresponding eigenvectors have
components
\begin{equation}
    \boxed{
    v_k^{(j)}
    =
    \sin\left(
        \frac{kj\pi}{n+1}
    \right),
    \qquad k=1,\ldots,n.
    }
\end{equation}

Since the angles $\theta_j$ are
distinct and cosine is strictly
decreasing on $(0,\pi)$, these
$n$ eigenvalues are distinct.

Thus, we have determined the
complete spectrum of $T_J$
without computing a determinant
or performing a numerical
eigenvalue calculation.

%```latex
We can now use these eigenvalues to determine
the optimal relaxation parameter for our
finite-difference system. Recall that the
spectral radius of a matrix is the largest
absolute value among its eigenvalues. Since
\[
    \mu_j=
    \cos\left(\frac{j\pi}{n+1}\right),
    \qquad j=1,\ldots,n,
\]
and the cosine function is strictly decreasing
on $(0,\pi)$, the eigenvalues of largest
absolute value occur at $j=1$ and $j=n$.
Indeed,
\[
    \mu_1=
    \cos\left(\frac{\pi}{n+1}\right),
\]
while
\[
\begin{aligned}
    \mu_n
    &=
    \cos\left(\frac{n\pi}{n+1}\right)\\
    &=
    \cos\left(\pi-\frac{\pi}{n+1}\right)\\
    &=
    -\cos\left(\frac{\pi}{n+1}\right).
\end{aligned}
\]
Consequently, the spectral radius of the
Jacobi iteration matrix is
\begin{equation}
    \boxed{
    \rho(T_J)=
    \cos\left(\frac{\pi}{n+1}\right).
    }
    \label{eq:jacobi_fd_spectral_radius}
\end{equation}

We are now ready to apply our theoretical
result. Recall that the optimal SOR
relaxation parameter is
\[
    \omega_{\mathrm{opt}}
    =
    \frac{2}{1+\sqrt{1-\rho(T_J)^2}}.
\]
Substituting the spectral radius obtained
above, we have
\[
\begin{aligned}
    \omega_{\mathrm{opt}}
    &=
    \frac{2}{
        1+\sqrt{
        1-\cos^2\left(\frac{\pi}{n+1}\right)
        }
    }\\[6pt]
    &=
    \frac{2}{
        1+\sqrt{
        \sin^2\left(\frac{\pi}{n+1}\right)
        }
    }\\[6pt]
    &=
    \frac{2}{
        1+\sin\left(\frac{\pi}{n+1}\right)
    }.
\end{aligned}
\]
Here, we used the identity
$1-\cos^2\theta=\sin^2\theta$ and the
fact that $\sin(\pi/(n+1))>0$.
Thus, for the standard one-dimensional
finite-difference matrix with $n$ interior
unknowns, we obtain
\begin{equation}
    \boxed{
    \omega_{\mathrm{opt}}
    =
    \frac{2}{
        1+\sin\left(\frac{\pi}{n+1}\right)
    }.
    }
    \label{eq:sor_optimal_fd}
\end{equation}
This is a remarkable result: the optimal
relaxation parameter depends only on the
number of interior grid points, and no
numerical eigenvalue computation is needed.
Moreover, as $n\to\infty$, we have
$\sin(\pi/(n+1))\to0$, and therefore
$\omega_{\mathrm{opt}}\to2$.

%```latex
\section{Numerical Experiments and Discussion}

We have established that the optimal relaxation
parameter for the standard one-dimensional
finite-difference matrix with $n$ interior
unknowns is

\begin{equation}
    \boxed{
    \omega_{\mathrm{opt}}
    =
    \frac{2}{
    1+\sin\left(\frac{\pi}{n+1}\right)
    }.
    }
    \label{eq:sor_numerical_optimal}
\end{equation}

Our next objective is to investigate how this
theoretical result behaves in numerical
computation. In particular, we wish to examine
how the choice of the relaxation parameter
affects the convergence of the SOR method
and whether the theoretically optimal value
agrees with our numerical observations.

\subsection{Experimental Setup}

For our numerical experiment, we consider
the boundary value problem

\[
    u''(x)=\sin x,
    \qquad 0<x<1,
\]

subject to the homogeneous Dirichlet
boundary conditions

\[
    u(0)=u(1)=0.
\]

This problem has the exact solution

\[
    u(x)=x\sin(1)-\sin x,
\]

which can be verified by direct
differentiation and substitution into
the boundary conditions.

We divide the interval $[0,1]$ into
$N=500$ equal subintervals, giving

\[
    h=\frac{1}{500}.
\]

There are therefore $n=N-1=499$
interior unknowns.

Using the central-difference
approximation, we obtain

\[
    \frac{
    u_{i-1}-2u_i+u_{i+1}
    }{h^2}
    =
    \sin(x_i).
\]

Multiplying by $-h^2$, we arrive at
the linear system

\[
    Au=b,
\]

where

\[
    A=\operatorname{tridiag}(-1,2,-1),
\]

and

\[
    b_i=-h^2\sin(x_i).
\]

This is precisely the class of symmetric
positive definite tridiagonal matrices
for which we derived the optimal SOR
relaxation parameter.

Substituting $n=499$ into
\eqref{eq:sor_numerical_optimal}, we obtain

\[
\begin{aligned}
    \omega_{\mathrm{opt}}
    &=
    \frac{2}{
    1+\sin\left(\frac{\pi}{500}\right)
    }\\[6pt]
    &\approx 1.9875.
\end{aligned}
\]

Thus, our theoretical analysis predicts
that the fastest asymptotic convergence
should occur near

\[
    \boxed{\omega_{\mathrm{opt}}\approx1.9875.}
\]

\subsection{Numerical Results}

To investigate this prediction, we
implemented the SOR method and repeated
the computation for several values
of the relaxation parameter $\omega$.

The stopping tolerance was set to
$10^{-6}$, and we recorded the number
of iterations required for convergence.

The results are summarized in
Table~\ref{tab:sor_numerical_results}.

\begin{table}[htbp]
\centering
\caption{SOR iteration counts for different
relaxation parameters with $N=500$.}
\label{tab:sor_numerical_results}
\begin{tabular}{cc}
\hline
$\omega$ & Iterations \\
\hline
1.9700 & 4082 \\
1.9750 & 3398 \\
1.9800 & 2680 \\
1.9850 & 1877 \\
1.9875 & 1320 \\
1.9900 & 1502 \\
1.9950 & 3002 \\
\hline
\end{tabular}
\end{table}

The numerical results are consistent
with our theoretical prediction.

As $\omega$ increases from $1.970$
toward the theoretically optimal value,
the number of iterations decreases
substantially.

For example, choosing $\omega=1.970$
requires $4082$ iterations, whereas
choosing $\omega\approx1.9875$
requires only $1320$ iterations.

However, increasing $\omega$ beyond
this value does not produce further
improvement. At $\omega=1.990$,
the iteration count increases to
$1502$, and at $\omega=1.995$,
it increases to $3002$.

This behavior illustrates the
mathematical phenomenon established
in our earlier analysis.

Before the optimal parameter,
increasing $\omega$ reduces the
dominant convergence factor. At the
optimal parameter, the spectral
radius reaches its minimum. Beyond
this point, the spectral radius
increases, and convergence becomes
progressively slower.

It is important to distinguish
between the theoretical and numerical
notions of optimality. Our formula
minimizes the spectral radius of the
SOR iteration matrix, which governs
the asymptotic rate of convergence.
The precise iteration count observed
in a finite computation can also
depend on the initial approximation,
the stopping criterion, and the
distribution of the initial error
among the eigenmodes.

Nevertheless, the close agreement
between the theoretical prediction
and the numerical results provides
a useful illustration of the
underlying theory.

\subsection{Implementation and Further Exploration}

The complete computational notebook
used for this exploration is publicly
available in the author's GitHub
repository:

\begin{center}
\url{https://github.com/biajid/Numerical-Linear-Algebra-/blob/main/Sauer_Problems/SOR_Exploration.ipynb}
\end{center}

Interested readers are encouraged
to examine the implementation,
reproduce the numerical experiments,
and investigate how the convergence
behavior changes with the number
of grid points, the relaxation
parameter, and the stopping tolerance.

The implementation was developed
primarily for mathematical exploration
and was optimized only cursorily.
Consequently, there is considerable
room for improving its computational
efficiency.

For example, readers may exploit
the tridiagonal structure of the
finite-difference matrix to avoid
unnecessary dense matrix operations,
investigate more efficient
implementations of the SOR iteration,
or compare the computational cost
with other iterative and direct
methods.

Such experiments provide an
opportunity to appreciate an
important distinction in numerical
linear algebra: a method may have
an excellent theoretical convergence
rate, but its practical efficiency
also depends on how carefully
the algorithm is implemented.

Our numerical experiment illustrates
how theoretical analysis can guide
the selection of an algorithmic
parameter, while computational
experimentation provides a means
of examining the consequences
of that choice in practice.


\section{Conclusion}

We began this note with one of the most fundamental
problems in scientific computing: solving a
system of linear equations,
\[
    Ax=b.
\]
Although direct methods such as Gaussian elimination
and LU or QR factorization provide systematic approaches
to this problem, the size and structure of many
systems arising in scientific computation make
iterative methods particularly attractive.

We first introduced the Jacobi and Gauss--Seidel
methods and examined how successive approximations
can converge to the exact solution. By studying
their iteration matrices, we established the
connection between convergence and the spectral
radius. This led us naturally to the successive
over-relaxation (SOR) method, in which a relaxation
parameter $\omega$ allows us to modify the
Gauss--Seidel updates in an attempt to accelerate
convergence.

However, introducing a relaxation parameter
immediately raises another question: how should
we choose it? An arbitrary choice of $\omega$
may improve convergence, produce little benefit,
or even cause the iteration to diverge. Thus,
the effectiveness of SOR depends not merely
on introducing relaxation, but on choosing
the relaxation parameter appropriately.

To investigate this question, we studied the
convergence properties of SOR and established
the importance of the interval
\[
    0<\omega<2.
\]
We then developed a relationship between the
eigenvalues of the Jacobi and SOR iteration
matrices for the class of symmetric positive
definite tridiagonal systems under consideration.
By examining the resulting quadratic equation,
its discriminant, and the behavior of its roots,
we obtained the optimal relaxation parameter
\[
    \boxed{
    \omega_{\mathrm{opt}}
    =
    \frac{2}{
        1+\sqrt{1-\rho(T_J)^2}
    }.
    }
\]

We subsequently applied this theoretical result
to the standard one-dimensional finite-difference
matrix. By solving the eigenvector recurrence
analytically, we determined the eigenvalues of
the corresponding Jacobi iteration matrix and
obtained the explicit expression
\[
    \boxed{
    \omega_{\mathrm{opt}}
    =
    \frac{2}{
        1+\sin\left(\frac{\pi}{n+1}\right)
    }.
    }
\]
This formula allows us to determine the optimal
relaxation parameter directly from the number
of interior grid points, without performing
a numerical eigenvalue calculation.

Finally, we examined the theory through a
numerical experiment involving a finite-difference
approximation to a boundary value problem.
The observed iteration counts agreed closely
with the theoretical prediction. In particular,
we found that increasing $\omega$ toward its
optimal value substantially reduced the number
of iterations, while increasing it beyond that
value made convergence slower.

This illustrates an important principle of
numerical computation: a seemingly small change
in an algorithmic parameter can have a substantial
effect on computational efficiency.

Finally, the historical importance of relaxation methods
should not be overlooked. Their development
was closely connected with the numerical solution
of partial differential equations, particularly
those arising in mathematical physics and
engineering. During the middle decades of the
twentieth century, when computing resources
were far more limited than they are today,
reducing the number of iterations required to
solve a large system could represent a major
computational improvement.

The development and analysis of relaxation
methods contributed significantly to the
progress of scientific computing through the
1950s, 1960s, and 1970s. Problems involving
heat conduction, fluid mechanics, electrostatics,
and other physical phenomena increasingly
required the numerical solution of large
systems of equations. In such settings,
understanding convergence and choosing
effective iterative parameters were not
merely theoretical concerns; they were
essential parts of making numerical
computation practical.

There is also a broader lesson in the
mathematical development presented here.
The ideas underlying stationary iterative
methods, relaxation, convergence analysis,
and optimal parameter selection emerged
through the contributions of numerous
mathematicians over several decades.
The theory did not appear fully developed
in a single paper or at a single moment.
Rather, it grew through successive
improvements, new mathematical insights,
and the increasing demands of scientific
computation.

For readers encountering these arguments
for the first time, some parts of the
derivation may appear difficult, particularly
the connection between the Jacobi and SOR
eigenvalues and the subsequent analysis
of the discriminant. This is entirely
understandable. Mathematics that required
the collective efforts of many researchers
over a period approaching half a century
should not necessarily be expected to
become completely transparent after
only fifty minutes of reading.

Indeed, understanding mathematics is
itself an iterative process. We begin
with an approximation to an idea, examine
where our understanding is incomplete,
and return to it with a better perspective.
Each attempt brings us a little closer
to the desired result.

Perhaps this is an appropriate place
to conclude a note on iterative methods.
Just as the SOR algorithm approaches
the solution through successive
approximations, our understanding of
mathematics often develops through
successive attempts.

The important thing is not to understand
everything immediately, but to continue
the iteration.


\begin{thebibliography}{9}

\bibitem{strang2006}
Gilbert Strang,
\textit{Linear Algebra and Its Applications},
4th ed.,
Brooks/Cole, 2006.

\bibitem{young1971}
David M. Young,
\textit{Iterative Solution of Large Linear Systems},
Academic Press, 1971.

\bibitem{varga2000}
Richard S. Varga,
\textit{Matrix Iterative Analysis},
2nd ed.,
Springer, 2000.

\end{thebibliography}

\end{document}

